\documentclass[10pt,a4paper]{amsart}
\usepackage[margin=19mm]{geometry}
\usepackage{amsmath,amssymb,mathtools}
\usepackage[scr=rsfs,cal=euler]{mathalfa}
\usepackage{xfrac}
\usepackage[unicode,hidelinks,bookmarks=false]{hyperref}
\newcommand{\ie}{i.e.\ }
\newcommand{\cf}{cf.\ }
\newcommand{\resp}{respectively\ }
\newcommand{\Subsection}{\subsection}
\providecommand{\nobreakdash}{\nobreak\hbox{-}}
\setlength{\emergencystretch}{2em}
\multlinegap0pt
\usepackage{amssymb}
\usepackage{braket}
\newtheorem{lem}{Lemma}
\title{A degeneration of the generalized Zwegers' $\mu$-function according to the Ramanujan difference equation}
\author{G. Shibukawa, S. Tsuchimi}
\date{Source: arXiv:2603.03693v2 (CC BY 4.0)}
\begin{document}
\maketitle

\noindent\emph{Source and licence.} A degeneration of the generalized Zwegers' $\mu$-function according to the Ramanujan difference equation. Version arXiv:2603.03693v2 (CC BY 4.0), distributed by arXiv under the Creative Commons Attribution 4.0 International licence, http://creativecommons.org/licenses/by/4.0/, as recorded in the arXiv licence metadata for this exact version and shown on the abstract page https://arxiv.org/abs/2603.03693v2. Full licence notice: \texttt{licenses/arxiv-2603.03693-CC-BY-4.0.txt}.\par\smallskip

\noindent\textit{Editorial scope.} Counted unit: the article's lem q-Fibonacci (source0.tex lines 1238--1245). The article's complete original proof of it is delivered with it (source0.tex lines 1246--1262, one \texttt{proof} environment of 17 lines). No mathematics is written, repaired or re-derived: the statement and the proof are the article's own lines, in the article's own notation, and no retained line is substituted or re-flowed.  Retained uncounted as the source-local context the proof needs: lines 299--302.  Nothing was omitted from any retained span, so the package carries no omission record.  No retained line contains a citation, so the wrapper carries no bibliography and no bibliography entry.  No retained line contains a figure environment, an image-inclusion command or drawing code, so no image is delivered and the package declares no asset dependency.  Editorial formatting only, in addition: an AMS \texttt{amsart} preamble that restates the article's own declarations and macros used by the retained lines, namely the environments \texttt{lem} on separate counters, together with the standard packages those lines need.  The retained environments print in this document as Lemma 1 in order of appearance, whereas the article prints the same items with the numbering of its own full document.  The complete native source of the article as distributed in the arXiv e-print for this exact version is bundled unchanged as \texttt{sources/195691/source0.tex}.
\begin{align}
\label{eq: q-Fibonacci}
F_n(t,q)=F_{n-1}(t,q)+tq^{n-2}F_{n-2}(t,q)
\end{align}

\begin{lem}
\label{lem: q-Fibonacci}
$(1)$\ If a sequence $F_n(t,q)$ satisfies \eqref{eq: q-Fibonacci}, then the sequence $(-t)^nq^\frac{n(n-1)}{2}F_{-n+1}\left(t,q^{-1}\right)$ also satisfies \eqref{eq: q-Fibonacci}. \\
$(2)$\ For all integers $n$, we have: 
\begin{align*}
S_n(t,q)=(-t)^{n-1}q^\frac{n(n-1)}{2}T_{-n}\left(t,q^{-1}\right),\quad T_{n-1}(t,q)=(-t)^{n}q^\frac{n(n-1)}{2}S_{-n+1}\left(t,q^{-1}\right). 
\end{align*}
\end{lem}

\begin{proof}
(1)\ It is clear. \\
(2)\ 
Since the sequences $S_n(t,q)$ and $T_{n-1}(t,q)$ satisfy \eqref{eq: q-Fibonacci}, the sequences 
\begin{align*}
S_n'(t,q):=(-t)^{n}q^\frac{n(n-1)}{2}S_{-n+1}\left(t,q^{-1}\right),\quad T_{n-1}'(t,q):=(-t)^{n-1}q^\frac{n(n-1)}{2}T_{-n}\left(t,q^{-1}\right)
\end{align*}
also satisfy \eqref{eq: q-Fibonacci} from $(1)$ of this Lemma. 
Moreover, since 
\begin{align*}
S_0'(t,q)=1=T_{-1}(t,q), \quad S_1'(t,q)=0=T_0(t,q),\quad T_{-1}'(t,q)=0=S_0(t,q), \quad T_{0}'(t,q)=1=S_1(t,q)
\end{align*}
we have
\begin{align*}
S_n'(t,q)=T_{n-1}(t,q),\quad T_{n-1}'(t,q)=S_n(t,q). 
\end{align*}
\end{proof}

\end{document}
